% TYPICAL COMMANDS FOR LECTURES

\renewcommand{\emph}[1]{{\color{blue}\textbf{#1}}}

\newcommand{\deutsch}[1]{{\color{blue}(#1)}}
\newcommand{\deutschertitel}[1]{{\tiny\deutsch{#1}}}

\newcommand{\mycite}[1]{``#1''}
\newcommand{\mycitebegin}{``}
\newcommand{\myciteend}{''}
\newcommand{\mytitlesource}[1]{{\tiny\normalfont\mbox{[#1]}}}
\newcommand{\mysource}[1]{\ifthenelse{\equal{#1}{}}{}{\phantom{.}~\hfill~\mytitlesource{#1}}}
\newcommand{\mychapter}[1]{, Chapter~#1}
\newcommand{\mypage}[1]{, p.~#1}
\newcommand{\mypages}[1]{, pp.~#1}

\newcommand{\todo}[1]{{\color{red}\textbf{[#1]}}}
\newcommand{\fodo}[1]{\todo{\footnote{\todo{#1}}}}
\newcommand{\todots}{\todo{\ldots}}

\newcommand{\textheightoftitle}{.125\textheight}
\newcommand{\textheightwithtitle}{.825\textheight}
\newcommand{\textheightwithouttitle}{.95\textheight}

\newcommand{\lessonslearned}[3]{
	\subsection{Summary}
	\begin{frame}{\insertsection{} -- \myframetitle}
		\begin{mycolumns}[animation=none]
			\begin{definition}{Lessons Learned}
				\begin{itemize}
					#1
				\end{itemize}
			\end{definition}
			\uncover<2->{\begin{note}{Further Reading}
				\small % references take space, can be a little smaller
				\begin{itemize}
					#2
				\end{itemize}
			\end{note}}
		\mynextcolumn
			\uncover<3->{\begin{example}{Practice}
				#3
			\end{example}}
		\end{mycolumns}
	\end{frame}
}

\newcommand{\faq}[3]{
	\subsection{FAQ}
	\begin{frame}{\myframetitle: \insertsubtitle{}}
		\small
		\begin{mycolumns}[t,columns=3,widths={34,34,34}]
			Part~\parta
			\begin{itemize}
				#1
			\end{itemize}
		\mynextcolumn
			Part~\partb
			\begin{itemize}
				#2
			\end{itemize}
		\mynextcolumn
			Part~\partc
			\begin{itemize}
				#3
			\end{itemize}
		\end{mycolumns}
	\end{frame}
}

% REFERENCES TO OTHER LECTURES

\newcommand{\reflecture}[1]{{Lecture~#1}}
\newcommand{\lectureintroduction}{\reflecture{1}}
\newcommand{\lectureruntime}{\reflecture{2}}
\newcommand{\lecturecloneandown}{\reflecture{3}}
\newcommand{\lecturemodeling}{\reflecture{4}}
\newcommand{\lecturefeatures}{\reflecture{5}} % formely \lectureembedded
\newcommand{\lecturemodules}{\reflecture{6}} % formerly \lectureapps
\newcommand{\lecturelanguages}{\reflecture{7}}
\newcommand{\lectureprocess}{\reflecture{8}}
\newcommand{\lectureinteractions}{\reflecture{9}}
\newcommand{\lectureanalyses}{\reflecture{10}}
\newcommand{\lecturetesting}{\reflecture{11}}
\newcommand{\lectureevonance}{\reflecture{12}}
\newcommand{\parta}{a}
\newcommand{\partb}{b}
\newcommand{\partc}{c}
% TODO how to handle different numberings for different universities?

% COMMANDS FOR PROPOSITIONAL FORMULAS AND MATHEMATICAL NOTATIONS

\newcommand{\sem}[1]{\ensuremath{\llbracket #1 \rrbracket}} % semantics brackets
\newcommand{\pand}{\wedge} % conjunction
\newcommand{\por}{\vee} % disjunction
\newcommand{\pnot}{\neg} % negation
\newcommand{\pequals}{\leftrightarrow} % biconditional
\newcommand{\npequals}{\nleftrightarrow} % exclusive disjunction
\newcommand{\mequals}{\Leftrightarrow} % equivalence (meta-level)
\newcommand{\pimplies}{\rightarrow} % conditional
\newcommand{\mimplies}{\Rightarrow} % implication (meta-level)
\newcommand{\defeq}{\vcentcolon=} % defining equals
\newcommand{\power}[1]{\mathcal{P}(#1)} % power set
\newcommand{\refslide}[1]{\hyperlink{#1}{(see Slide \autoref{#1})}} % link to slide
\newcommand{\fs}[1]{\ensuremath{{\color{green}#1}}} % selected feature
\newcommand{\fd}[1]{\ensuremath{{\color{red}#1}}} % deselected feature
\newcommand{\clause}[1]{\ensuremath{{\color{red}#1}}} % clause in a CNF
\newcommand{\cfg}[3][0]{\ensuremath({{\only<#1>{\color{green}}\{#2\}}, {\only<#1>{\color{red}}\{#3\}})}} % configuration

\usepackage{mathtools} % required for absolute value in modeling lecture
\DeclarePairedDelimiter\abs{\lvert}{\rvert} % absolute value

% COMMANDS TO INCLUDE XKCDs

\newcommand{\xkcd}[2]{\pic[#2]{#1}}
\newcommand{\xkcdframe}[1]{
	\begin{frame}
		\centering%
		\xkcd{#1}{height=80mm}
	\end{frame}
}
\newcommand{\widexkcdframe}[1]{
	\begin{frame}
		\xkcd{#1}{width=\linewidth}
	\end{frame}
}

% COMMANDS TO SIMPLIFY USAGE OF THE PROJECT CARTOON

\newcommand{\projectcartoonwidth}{.19} % default value for 5 tiles
\newcommand{\projectcartoon}[2]{\projectcartoonimage{sepia/cell_#1}{#2}}
\newcommand{\hprojectcartoon}[2]{\projectcartoonimage{alpha/cell_#1}{#2}}
\newcommand{\projectcartoonimage}[2]{
	\begin{minipage}[t]{\projectcartoonwidth\linewidth}
		\centering
		\href{http://web.archive.org/web/20191029221320/http://www.projectcartoon.com/}{\includegraphics[width=\linewidth]{#1}}\\#2
	\end{minipage}\hfill
}
\newcommand{\waterfallcartoon}{
	\hprojectcartoon{02}{Analysis}%
	\hprojectcartoon{03}{Design}%
	\hprojectcartoon{04}{Implementation}%
	\hprojectcartoon{05}{Testing}%
	\hprojectcartoon{10}{Maintenance}%
}
%\uncover<1->{\hprojectcartoon{01}{how the customer explained it}} % requirements
%\uncover<1->{\hprojectcartoon{02}{how the project leader understood it}} % modeling
%\uncover<1->{\hprojectcartoon{03}{how the analyst designed it}} % architecture and design
%\uncover<1->{\hprojectcartoon{04}{how the programmer implemented it}} % implementation
%\uncover<1->{\hprojectcartoon{05}{what the beta testers received}} % testing
%\uncover<1->{\hprojectcartoon{06}{how the business consultant described it}} % process
%\uncover<1->{\hprojectcartoon{07}{how the project was documented}} % documentation?
%\uncover<1->{\hprojectcartoon{08}{what operations installed}} % devops/continuous integration?
%\uncover<1->{\hprojectcartoon{09}{how the customer was billed}} % management (pricing)
%\uncover<1->{\hprojectcartoon{10}{how it was supported}} % maintenance
%\uncover<1->{\hprojectcartoon{11}{what marketing advertised}} % product lines. marketing? reuse?
%\uncover<1->{\hprojectcartoon{12}{when it was delivered}} % configuration management. continous delivery?
%\uncover<1->{\hprojectcartoon{13}{what the customer really needed}} % customer / SE II
%\uncover<1->{\hprojectcartoon{14}{what the digg effect can do to your site}} % ???
%\uncover<1->{\hprojectcartoon{15}{the disaster recover plan}} % ???
%\uncover<1->{\hprojectcartoon{16}{the open source version}} % open source (licensing)
%\uncover<1->{\hprojectcartoon{17}{how it performed under load}} % compilation and static analyses. quality assurance? performance?
%\uncover<1->{\hprojectcartoon{18}{how patches were applied}} % evolution

% COMMANDS BASED ON TIKZ

\newcommand{\mglass}[1]{
	\scalebox{#1}{
		\begin{tikzpicture}[overlay]
			\node[circle,draw=gray,line width=1,fill=blue,fill opacity=.2,inner sep=2mm] {} edge[draw=gray,line width=1,double distance=1pt,shorten <= -.2pt] (.33,-.33);
		\end{tikzpicture}
	}
}

\def\checkmark{\tikz\fill[scale=0.4](0,.35) -- (.25,0) -- (1,.7) -- (.25,.15) -- cycle;}

% OTHER COMMANDS

\newcommand{\ssat}{\texttt{\#}SAT}